\documentclass[10pt,a4paper]{amsart}
\usepackage[margin=19mm]{geometry}
\usepackage{amsmath,amssymb,mathtools}
\usepackage[scr=rsfs,cal=euler]{mathalfa}
\usepackage{xfrac}
\usepackage[unicode,hidelinks,bookmarks=false]{hyperref}
\newcommand{\ie}{i.e.\ }
\newcommand{\cf}{cf.\ }
\newcommand{\resp}{respectively\ }
\newcommand{\Subsection}{\subsection}
\providecommand{\nobreakdash}{\nobreak\hbox{-}}
\setlength{\emergencystretch}{2em}
\multlinegap0pt
\usepackage{amssymb}
\usepackage{mathtools}
\usepackage{tikz}
\newtheorem{problem}{Problem}
\newtheorem{theorem}{Theorem}
\newcommand{\NP}{\mathsf{NP}}
\newcommand{\UG}{\operatorname{UG}}
\newcommand{\cB}{\mathcal{B}}
\title{The Complexity of Mixed Arc-Disjoint Spanning Subdigraphs with Antistrong Connectivity}
\author{Jiangdong Ai, Gregory Gutin, Hui Lei, Yongtang Shi}
\date{Source: arXiv:2608.00115v1 (CC BY 4.0)}
\begin{document}
\maketitle

\noindent\emph{Source and licence.} The Complexity of Mixed Arc-Disjoint Spanning Subdigraphs with Antistrong Connectivity. Version arXiv:2608.00115v1 (CC BY 4.0), distributed by arXiv under the Creative Commons Attribution 4.0 International licence, http://creativecommons.org/licenses/by/4.0/, as recorded in the arXiv licence metadata for this exact version and shown on the abstract page https://arxiv.org/abs/2608.00115v1. Full licence notice: \texttt{licenses/arxiv-2608.00115-CC-BY-4.0.txt}.\par\smallskip

\noindent\textit{Editorial scope.} Counted unit: the article's theorem thm:main-2EC (source0.tex lines 188--193). The article's complete original proof of it is delivered with it (source0.tex lines 500--626, one \texttt{proof} environment of 127 lines, which the article places away from the statement). No mathematics is written, repaired or re-derived: the statement and the proof are the article's own lines, in the article's own notation, and no retained line is substituted or re-flowed.  Retained uncounted as the source-local context the proof needs: lines 168--173; lines 451--453; lines 458--461.  Nothing was omitted from any retained span, so the package carries no omission record.  No retained line contains a citation, so the wrapper carries no bibliography and no bibliography entry.  No retained line contains a figure environment, an image-inclusion command or drawing code, so no image is delivered and the package declares no asset dependency.  Editorial formatting only, in addition: an AMS \texttt{amsart} preamble that restates the article's own declarations and macros used by the retained lines, namely the environments \texttt{problem}, \texttt{theorem} on separate counters, together with the standard packages those lines need.  The retained environments print in this document as Problem 1, Theorem 1 in order of appearance, whereas the article prints the same items with the numbering of its own full document.  The complete native source of the article as distributed in the arXiv e-print for this exact version is bundled unchanged as \texttt{sources/206518/source0.tex}.
\begin{problem}[Antistrong--2-edge-connected spanning packing]
\label{prob:AS-2EC}
Can we decide in polynomial time whether a digraph $D$ contains arc-disjoint spanning
subdigraphs $D_1,D_2$ such that $D_1$ is antistrong and
$\UG(D_2)$ is 2-edge-connected?
\end{problem}

\begin{equation}\label{eq:underlying-core}
   \UG(R)\cong K_7-02,
\end{equation}

\begin{problem}
\label{prob:PCBHAM}
The input is a planar cubic bipartite graph $G$. Decide whether $G$ contains a Hamiltonian path.
\end{problem}

\begin{theorem}\label{thm:main-2EC}
It is NP-complete to decide whether a digraph $D$ contains
arc-disjoint spanning subdigraphs $D_1,D_2$ such that $D_1$ is antistrong and $\UG(D_2)$ is 2-edge-connected. It remains NP-complete even for oriented digraphs $D$ that are strong and antistrong, whose underlying graphs $\UG(D)$ are
3-vertex-connected, and in which all but at most two vertices $v\in V(D)$ satisfy
$d^+(v)\leq4$ and $d^-(v)\leq4$.
\end{theorem}

\begin{proof}[Proof of Theorem~\ref{thm:main-2EC}]
The problem belongs to $\NP$: antistrong connectivity is checked by connectivity of the bipartite representation, and 2-edge-connectivity of the underlying graph is checked by a bridge
search.

We prove NP-hardness by a polynomial reduction from
Problem~\ref{prob:PCBHAM}.  Let $G=(U,W;E)$ be a connected instance,
and let $D(G)$ be the oriented graph constructed above.  We claim that $G$ has a Hamiltonian path if and only if $D(G)$ is a yes-instance of the decision problem described in
Problem~\ref{prob:AS-2EC}.
Suppose first that $P$ is a Hamiltonian path of $G$, and let $S=E\setminus E(P)$.  Define two spanning subdigraphs of $D(G)$ by
\begin{align*}
 A(D_{\rm anti})={}&F_7
   \cup\{1x_u:u\in U\}
   \cup\{x_w1:w\in W\}
   \cup\{x_{u_*}2\}\\
   &\cup\{x_ux_w:uw\in E(P)\},\\
 A(D_2)={}&C_7
   \cup\{0x_u:u\in U\}
   \cup\{x_w0:w\in W\}
   \cup\{x_ux_w:uw\in S\}.
\end{align*}
The two arc sets are disjoint.

Consider in $\cB(D_{\rm anti})$ the set $L=\{x_u':u\in U\}\cup\{x_w'':w\in W\}$.
The variable arcs corresponding to $E(P)$ induce on $L$ a graph
isomorphic to the Hamiltonian path $P$.  The connector is represented
by the edge $x_{u_*}'2''$ and therefore joins this path to the
connected graph $\cB(R[F_7])$.  For every $u\in U$, the spoke $1x_u$
attaches the remaining clone $x_u''$ to the core, and for every
$w\in W$, the spoke $x_w1$ attaches the remaining clone $x_w'$ to the
core.  Thus $\cB(D_{\rm anti})$ is connected, and
$D_{\rm anti}$ is antistrong.

Since $G$ is cubic, every internal vertex of $P$ is incident with one
edge of $S$, and each end of $P$ is incident with two edges of $S$.
In particular,
\begin{equation}\label{eq:S-meets-forward}
   d_S(v)\ge 1\qquad\text{for every }v\in V(G).
\end{equation}
The graph $\UG(D_2)$ contains the core cycle $\UG(R[C_7])$.  Every terminal $x_v$ has a spoke to $0$.  For each $uw\in S$, the variable
edge $x_ux_w$ and the two corresponding spokes form the triangle $0x_ux_w0$.
Every variable edge lies on such a triangle, and by
\eqref{eq:S-meets-forward}, every spoke lies on at least one such
triangle.  The core edges lie on the core $7$-cycle.  Hence
$\UG(D_2)$ is connected and every one of its edges lies on a cycle.
It is therefore 2-edge-connected.

Conversely, suppose that $D(G)$ contains arc-disjoint spanning
subdigraphs $D_{\rm anti}$ and $D_2$, where $D_{\rm anti}$ is
antistrong and $\UG(D_2)$ is 2-edge-connected.  Let
\begin{align*}
 T&=\{uw\in E:x_ux_w\in A(D_{\rm anti})\},\\
 S&=\{uw\in E:x_ux_w\in A(D_2)\}.
\end{align*}
Clearly $S\cap T=\emptyset$.

We first show that $(V(G),T)$ is connected. In the full bipartite representation $\cB(D(G))$, let
\[
   L=\{x_u':u\in U\}\cup\{x_w'':w\in W\}.
\]
By construction,
\begin{equation}\label{eq:unique-connector-cut}
   \delta_{\cB(D(G))}(L)=\{x_{u_*}'2''\}.
\end{equation}
Indeed, every variable edge has both endpoints in $L$, every spoke is incident with one of the other terminal clones, and the connector is the only remaining edge incident with $L$.

Since $\cB(D_{\rm anti})$ is connected,
\eqref{eq:unique-connector-cut} forces the connector to belong to $D_{\rm anti}$.  Moreover, $\cB(D_{\rm anti})[L]$ must be connected.  Otherwise, a component not containing $x_{u_*}'$ would
have no edge to the rest of $\cB(D_{\rm anti})$. The only ambient edges with both endpoints in $L$ are the variable edges. Consequently, the graph $(V(G),T)$ is connected and spanning.

We next show that
\begin{equation}\label{eq:S-positive-oriented}
   d_S(v)\ge 1\qquad\text{for every }v\in V(G).
\end{equation}
For $u\in U$, the vertex $x_u''$ of the bipartite representation is
incident in the whole ambient graph only with the two spoke edges $0'x_u''$ and $1'x_u''$.  Connectivity of $\cB(D_{\rm anti})$ forces at least one of the corresponding spokes to
belong to $D_{\rm anti}$.  Similarly, for $w\in W$, the clone $x_w'$ is incident only with the two spoke edges $x_w'0''$ and $x_w'1''$, so again $D_{\rm anti}$ uses at least one spoke at $x_w$.

By arc-disjointness, $D_2$ therefore uses at most one of the two spokes at every terminal. Apart from its two spokes, every terminal is incident only with variable arcs, except that $x_{u_*}$ is also incident with the connector. The connector is unavailable to $D_2$, because it has already been forced into $D_{\rm anti}$.

Since a 2-edge-connected graph has minimum degree at least two, every terminal must therefore be incident in $D_2$ with at least one variable edge. This proves
\eqref{eq:S-positive-oriented}.

Finally, cubicity and arc-disjointness give $d_T(v)+d_S(v)\le 3$ for every $v\in V(G)$.
Together with \eqref{eq:S-positive-oriented}, this yields $d_T(v)\le 2$ for every $v$. The graph $(V(G),T)$ is connected and spanning, so it is a Hamiltonian path or a Hamiltonian cycle.
In the latter case, deleting one edge produces a Hamiltonian path. Hence $G$ has a Hamiltonian path, completing the equivalence.

It remains to verify the additional promises. The directed cycle $R[C_7]$ makes the core strong. The core reaches every $x_u$ directly through $0x_u$, and reaches every $x_w$ through a two-arc
path $0x_ux_w$, where $u$ is any neighbour of $w$.  Conversely, every $x_w$ reaches the core through $x_w0$, and every $x_u$ reaches the core through
$x_ux_w0$, where $w$ is any neighbour of $u$. Thus $D(G)$ is strong.

The graph $\cB(D(G))$ is connected as well. Indeed, use the connected core $\cB(R[F_7])$, all spokes incident with $1$, the connector, and all variable arcs. Since $G$ is connected, the variable edges connect all vertices of $L$, the connector joins $L$ to the core, and the spokes attach the remaining terminal clones. Hence $D(G)$ is antistrong.

We next prove that $\UG(D(G))$ is 3-vertex-connected. Let $H=\UG(D(G))$.
By \eqref{eq:underlying-core}, the subgraph of $H$ induced by the core vertices is isomorphic to $K_7-02$.

Let $Z\subseteq V(H)$ with $|Z|\le 2$. If at least one of the vertices $0,1$ is not in $Z$, then every remaining terminal is adjacent to a remaining vertex in $\{0,1\}$. Moreover, the remaining core is connected, since deleting at most two vertices from $K_7-02$ leaves a connected graph. Hence $H-Z$ is connected.

The only remaining case is
$Z=\{0,1\}$. The terminal vertices then induce a copy of the connected graph $G$, the remaining core is connected, and the connector
$x_{u_*}2$ joins the terminal copy to the core. Thus $H-\{0,1\}$ is connected. We conclude that $H$ is 3-vertex-connected. In particular, $H$ is also 2-edge-connected.

Finally, only the core vertices $0$ and $1$ can have in-degree or out-degree growing with $|V(G)|$. For the terminals,
\[
\begin{aligned}
 \bigl(d_{D(G)}^+(x_u),d_{D(G)}^-(x_u)\bigr)
   &=(3,2)
   &&\text{for }u\in U\setminus\{u_*\},\\
 \bigl(d_{D(G)}^+(x_{u_*}),d_{D(G)}^-(x_{u_*})\bigr)
   &=(4,2),\\
 \bigl(d_{D(G)}^+(x_w),d_{D(G)}^-(x_w)\bigr)
   &=(2,3)
   &&\text{for }w\in W.
\end{aligned}
\]
For the remaining core vertices, direct inspection gives
\[
\begin{array}{c|ccccc}
 v & 2 & 3 & 4 & 5 & 6\\ \hline
 d_{D(G)}^+(v) & 3 & 4 & 3 & 3 & 2\\
 d_{D(G)}^-(v) & 3 & 2 & 3 & 3 & 4
\end{array}
\]
where the connector is included in the in-degree of $2$. Consequently, every vertex other than $0$ and $1$ has both in-degree and out-degree at most four.

The fixed no-instance $D_0$ also satisfies all these promises, as verified above.  Therefore the reduction proves NP-hardness under
the restrictions stated in Theorem~\ref{thm:main-2EC}.  Since the problem belongs to $\NP$, the proof is complete.
\end{proof}

\end{document}
